\section*{Chapter \ref{ch:iv1:optimal-stopping} --- Optimal Stopping}
\markboth{\protect\omTocMark{}{Answers: Optimal Stopping}}{\protect\omTocMark{}{Answers: Optimal Stopping}}

\begin{solution}{iq:iv1:optimal-stopping:1}
With one re-roll, keep 4, 5 or 6, and the game is worth $4.25$. With two re-rolls, keep 5 or 6 at first, and the game is worth $14/3 \approx 4.67$.

Work backwards\index{backward induction}. The last roll is worth $3.5$, so with one re-roll left you keep anything above $3.5$. The value is then $\tfrac12 \cdot 5 + \tfrac12 \cdot 3.5 = 4.25$. With two re-rolls left, you compare the roll with $4.25$, so you keep 5 and 6 only. The value is $\tfrac13 \cdot 5.5 + \tfrac23 \cdot 4.25 = 14/3$. With many re-rolls the value climbs towards 6.

\iqwrong{Using the same threshold at every stage. The threshold is the value of what you would be giving up, and that value is higher when you have more rolls left.}
\end{solution}

\begin{solution}{iq:iv1:optimal-stopping:2}
The values are $5/8$, $89/128 \approx 0.695$ and $24305/32768 \approx 0.742$. At each stage the threshold is the value of the game that you would move to: $\tfrac12$, then $\tfrac58$, then $\tfrac{89}{128}$.

If the game is worth $v$ after a rejection, you keep any draw above $v$. So the new value is
\[
E[\max(U, v)] = v \cdot v + \int_v^1 x\,dx = \frac{1 + v^2}{2}.
\]
Start from $v = \tfrac12$, the value of a last draw that you must keep.

The recursion has its only fixed point at $v = 1$, so with many draws the value tends to the top of the range, as it should.

\iqwrong{Giving $2/3$ for one rejection, which is the mean of the larger of two draws. That is what you get with hindsight. Here you have to decide before you see the second draw.}
\end{solution}

\enlargethispage{-3\baselineskip}%
\begin{solution}{iq:iv1:optimal-stopping:3}
Keep 51 or higher, and the game is worth about $\$63.13$.

If you hold a ticket with number $x$, the second draw is uniform on the other 99, so its mean is $(5050 - x)/99$. You keep when
\[
x \ge \frac{5050 - x}{99} \iff 100x \ge 5050 \iff x \ge 50.5.
\]
The threshold is exactly the plain mean, $50.5$. Removing a low ticket raises the mean of the rest, removing a high one lowers it, and the two effects balance there.

The value is $\frac{1}{100}\bigl(\sum_{x = 51}^{100} x + \sum_{x = 1}^{50} \frac{5050 - x}{99}\bigr) = \frac{1}{100}(3775 + 2537.6)$. The continuous version, with a threshold of $\tfrac12$, is worth $0.625$, which is close to this on a scale of 100.

\iqwrong{Thinking that you should switch only for tickets far below the middle, because of ``the risk of losing what you have''. You are paid a mean, and a risk-neutral player should use the threshold at the mean.}
\end{solution}

\begin{solution}{iq:iv1:optimal-stopping:4}
No rule does better than $1/2$, and every rule achieves exactly $1/2$.

The fraction of red among the remaining cards is a martingale\index{martingale}. If $r$ of the $r + b$ remaining cards are red, the expected fraction after the next card is
\[
\frac{r}{r + b} \cdot \frac{r - 1}{r + b - 1} + \frac{b}{r + b} \cdot \frac{r}{r + b - 1} = \frac{r}{r + b}.
\]
Your win chance when you stop is the current fraction. Optional stopping\index{optional stopping} says that the expected value at any stopping rule is the starting value, $\tfrac12$.

A shorter proof: given what you have seen, the next card and the very last card are exchangeable. So stopping and betting on the next card has the same chance as betting on the last card of the deck, and that card is red with chance $\tfrac12$, whatever your rule.

\iqwrong{Waiting until more black cards have gone by, because then the deck is richer in red. It does get richer when that happens, but it might never happen, and on average those paths balance. A backward induction shows that every state is exactly indifferent.}
\iqwhy{This is the simplest example of ``no timing rule beats a fair game''. Anyone who shows you a rule that seems to get an edge from timing a fair process has made an error, or they are selecting their sample.}
\end{solution}

\begin{solution}{iq:iv1:optimal-stopping:5}
For $n = 3$: reject the first, and then take the first candidate better than all before, for a win chance of $1/2$. For $n = 10$: reject the first 3, for a win chance of about $0.399$. For large $n$: reject a fraction $1/e$, for a win chance tending to $1/e \approx 0.368$.

A good rule rejects the first $s$ and then takes the next candidate who beats all before\index{secretary problem}. You win if the best overall is at position $k > s$, with chance $\tfrac1n$ for each $k$, and if the best of the first $k - 1$ lies in the first $s$, with chance $\tfrac{s}{k - 1}$. So
\[
P(\text{win}) = \frac{s}{n} \sum_{k = s + 1}^{n} \frac{1}{k - 1}.
\]
For $n = 3$ and $s = 1$, this is $\tfrac13(1 + \tfrac12) = \tfrac12$. Taking anyone at a fixed position wins with chance only $\tfrac13$. For large $n$, with $s = xn$, the sum is about $\ln(1/x)$, so you maximise $x\ln(1/x)$, which peaks at $x = 1/e$ with value $1/e$.

The threshold form need not be assumed. At position $k$, a candidate who beats all before is the overall best with chance $k/n$, and that rises with $k$, while the value of going on falls. A backward induction on $k$ gives the same answer, $0.399$ for $n = 10$.

\iqwrong{Assuming the win chance falls to 0 as $n$ grows, because the best is ever harder to find. The win chance tends to a constant, however many candidates there are.}
\end{solution}

\begin{solution}{iq:iv1:optimal-stopping:6}
Stop at the first number above $0.8$. You make 5 draws on average, and the game is worth $0.8$. In general the threshold is $1 - \sqrt{2c}$.

Let $v$ be the value of the game before a draw that you must pay for\index{reservation price}. The future looks the same after any rejection, so a draw $x$ is kept when $x \ge v$. The value equation is $v = E[\max(U, v)] - c$, which says
\[
\int_v^1 (x - v)\,dx = c, \quad\text{so}\quad \frac{(1 - v)^2}{2} = c.
\]
With $c = 0.02$, you get $1 - v = 0.2$. The number of draws is geometric with success chance $0.2$, so its mean is 5. As a second route: you end up with a number averaging $0.9$, and you pay $5 \cdot 0.02$, which leaves $0.8$.

The threshold falls as the cost rises. At $c = \tfrac12$ the threshold is 0 and the game is worth nothing, because the first draw is worth just its cost. At $c = 0$ the threshold is 1, and you go on forever.

\iqwrong{Lowering the threshold as the costs pile up. The money already spent is gone, and the future looks the same after every draw, so there is one threshold for the whole game.}
\iqwhy{Going on is an option with a price. This same equation appears when you decide how long to search for a better price, or when to stop collecting data.}
\end{solution}

\begin{solution}{iq:iv1:optimal-stopping:7}
Keep 15 or higher. The game is worth $91/6 \approx \$15.17$.

Let $w$ be the value of the game when you are about to make a free roll. A re-roll therefore gives you $w - 1$, so you keep any number above $w - 1$. A threshold rule ``keep at least $t$'' has a kept number averaging $(t + 20)/2$, and $20/(21 - t)$ rolls on average, of which all but the first cost \$1. Its value is
\[
\frac{t + 20}{2} - \Bigl(\frac{20}{21 - t} - 1\Bigr).
\]
This gives $15.14$ at $t = 14$, $15.17$ at $t = 15$, and $15$ at $t = 16$. So $t = 15$ is best. It agrees with the condition: $w - 1 = 14.17$, so you keep what beats $14.17$, which is 15 or more. Value iteration on $w = E[\max(k, w - 1)]$ reaches the same value.

Limits: with a free re-roll the value would be 20. With a very high cost you keep the first roll, for a value of $10.5$.

\iqwrong{Comparing the roll with $w$ itself, which gives a threshold of 16 and a value of only 15. A re-roll costs \$1, so the number to beat is $w - 1$.}
\end{solution}

\begin{solution}{iq:iv1:optimal-stopping:8}
Accept anything at or above $0.5$. It takes 2 days on average, and the block is worth $0.625$ just before the first offer.

Let $v$ be the value before an offer arrives\index{discounted search}. A rejected offer gives you $\delta v$, with $\delta = 0.8$, so you accept when the offer beats $\delta v$. Then
\[
v = E[\max(U, \delta v)] = \frac{1 + (\delta v)^2}{2}.
\]
Solving this gives $v = \frac{1 - \sqrt{1 - \delta^2}}{\delta^2} = \frac{1 - 0.6}{0.64} = 0.625$, and the threshold is $\delta v = 0.5$. Half of the offers are accepted, so the wait is geometric, with mean 2 days. A simulation with this rule gives the same value.

The limits make sense. At $\delta = 0$ you take any offer, which is worth $\tfrac12$. At $\delta = 1$ you wait for an offer of nearly 1, and the value tends to 1.

\iqwrong{Waiting for $0.625$, the value of the block, a rule worth only about $0.609$. What you give up by refusing is that value a day later, so the threshold is $\delta v$, not $v$. That the threshold equals the mean offer, $0.5$, is a coincidence of $\delta = 0.8$.}
\end{solution}

\begin{solution}{iq:iv1:optimal-stopping:9}
Keep rolling below 15, and stop at 15 or more (you are indifferent at exactly 15). The game is worth about $\$6.15$.

Holding a total of $S$, one more roll gives an expected
\[
\frac16 \sum_{k = 1}^{5} (S + k) = \frac{5S}{6} + \frac{15}{6}.
\]
This beats $S$ exactly when $S < 15$. A rule that compares stopping with one more step and then stopping is often called one-step look-ahead\index{one-step look-ahead}. It is optimal here, because once it says ``stop'', it continues to do so at every higher total, and totals only go up.

For the value, use backward induction from totals of 15 and up, where the value is the total. For $S < 15$ the value is $\frac16\sum_{k = 1}^5 V(S + k)$. Working down gives $V(0) \approx 6.15$. A simulation of the rule matches.

Note how much the risk of a zero costs. Even if you follow the rule, you fail to reach 15 about 62\% of the time, and all of the value comes from the paths that do.

\iqwrong{Stopping when the total reaches $3.5 \cdot 5 = 17.5$, or after a set number of rolls. The rule depends on what you hold, because what you risk is the total, whereas the gain is a fixed $2.5$ on average, so the risk and reward come into balance at 15.}
\end{solution}

\begin{solution}{iq:iv1:optimal-stopping:10}
With negative payoffs allowed, every rule is worth 0. With the floor at 0, going to the end is optimal: the game is worth $3/4$.

Heads minus tails, $S_k$, is a fair random walk, hence a martingale\index{optional stopping}. For a bounded stopping time, $E[S_\tau] = S_0 = 0$. So all rules are worth the same: stopping when ahead brings gains offset by the losses of the paths that had to run to the end.

With a floor, the payoff is $\max(S_k, 0)$, a convex function of a martingale, hence a submartingale\index{submartingale}. Its expectation can only rise with time, so waiting until the end is optimal. Then
\[
E[S_4^+] = \frac{1 \cdot 4 + 4 \cdot 2}{16} = \frac34,
\]
from $S_4 = 4$ with chance $\tfrac{1}{16}$, and $S_4 = 2$ with chance $\tfrac{4}{16}$. For large $n$, this grows like $\sqrt{n/2\pi}$.

\iqwrong{Stopping at a gain of 1 on the first flip to lock in a profit, with the floor in place. Stopping pays 1, and going on pays $5/4$, because the floor protects you from the fall. (Stopping at 2 after two heads happens to cost nothing, since the walk can no longer go below 0.)}
\iqwhy{This is why it is not optimal to exercise an American call on a stock with no dividends and zero rates early. The call's payoff is convex in a martingale, so early exercise only gives away the option's remaining value.}
\end{solution}

\begin{solution}{iq:iv1:optimal-stopping:11}
The report is worth at most 10. A perfect report is worth at most 20.

Without a report\index{value of information}, you trade, with an expected $0.5 \cdot 100 - 40 = 10$.

With the report, ``yes'' and ``no'' are equally likely, by symmetry. After ``yes'', the chance of a deal is $0.8$, so trading is worth $40$. After ``no'', the chance is $0.2$, so trading is worth $-20$, and you pass, which gives 0. So the value is $\tfrac12 \cdot 40 + \tfrac12 \cdot 0 = 20$, which is 10 above the original.

With a perfect report you trade only when the deal will close, and you make 60. So it is worth $\tfrac12 \cdot 60 = 30$, which is 20 more.

The value comes only from changing a decision. A report that is right 60\% of the time is worth nothing here: after ``no'' you would have a chance of $0.4$, which is worth 0, so it saves you nothing. That is an important check: accuracy above $\tfrac12$ does not mean that information has value.

\iqwrong{Paying for the higher expected profit after ``yes'', $40 - 10 = 30$. You do not get ``yes'' every time, and you must commit to the price before you know what the report says.}
\end{solution}

\begin{solution}{iq:iv1:optimal-stopping:12}
Flip coin B first, then stay with B after a head and switch to A after a tail. The game is worth $13/12$, against 1 for coin A twice.

Coin B has the same chance on the first flip, $\tfrac12$. But it also tells you something\index{explore versus exploit}. After a head, the chance for B is $\tfrac23$ by Laplace's rule; after a tail it is $\tfrac13$, so you go to A. So the value is
\[
\tfrac12 + \tfrac12 \cdot \tfrac23 + \tfrac12 \cdot \tfrac12 = \tfrac{13}{12}.
\]
The extra $1/12$ is the value of learning. With only one flip left there is nothing to learn from, and both coins are worth $\tfrac12$. With three flips, the value is $5/3$, against $3/2$ for coin A, so an edge of $1/6$. The value of exploring grows with the time left to use what you learn.

\iqwrong{Saying you are indifferent because the two coins have the same mean. The unknown coin has the same mean, but it has variance in what you will know next, and you can exploit that.}
\iqwhy{Given two strategies with equal expected return, allocate the first capital to the one you know less about, as long as you will have time to act on what you learn.}
\end{solution}

\begin{solution}{iq:iv1:optimal-stopping:13}
Stop once the total is at least $\sqrt{2} - 1 \approx 0.414$. The game is worth $(\sqrt{2} - 1)e^{\sqrt{2} - 1} \approx 0.627$ at the start.

With a total of $s$, one more draw and then stopping has an expected payoff of
\[
\int_0^{1 - s} (s + u)\,du = s(1 - s) + \frac{(1 - s)^2}{2} = \frac{1 - s^2}{2}.
\]
This beats $s$ exactly when $s^2 + 2s - 1 < 0$, so when $s < \sqrt{2} - 1$\index{one-step look-ahead}. Once you exceed this, the total only goes up, so the rule stays ``stop''.

For the value at $s$ below the threshold $t$, the value function satisfies $V(s) = \int_s^1 V(x)\,dx$, where $V(x) = x$ above $t$. Differentiating gives $V' = -V$, so $V(s) = V(t)e^{t - s}$, and $V(t) = t$. So $V(0) = te^t$. A simulation of the rule gives the same, and scanning other thresholds confirms the maximum at $t$.

\iqwrong{Stopping once the total is above $\tfrac12$, by symmetry. There is no symmetry, because going over costs you everything you have, and at $\tfrac12$ that is a lot for a low chance of a modest gain.}
\end{solution}

\begin{solution}{iq:iv1:optimal-stopping:14}
The deck of 2 and 2 is worth $2/3$, and the full deck is worth about \$2.62.

Let $V(r, b)$ be the value with $r$ red and $b$ black left. Then $V(0, b) = 0$, $V(r, 0) = r$ and\index{backward induction}
\[
V(r, b) = \max\Bigl(0,\ \frac{r}{r + b}\bigl(1 + V(r - 1, b)\bigr) + \frac{b}{r + b}\bigl(V(r, b - 1) - 1\bigr)\Bigr).
\]
This is a $27 \times 27$ table, and you fill it in from the small corner. You stop when the second term is not positive.

By hand: $V(1, 1) = \tfrac12$. $V(1, 2)$ has continuation value $\tfrac13(1 + 0) + \tfrac23(\tfrac12 - 1) = 0$. And $V(2, 1) = \tfrac23(1 + \tfrac12) + \tfrac13(2 - 1) = \tfrac43$. So
\[
V(2, 2) = \tfrac12\bigl(1 + V(1, 2)\bigr) + \tfrac12\bigl(V(2, 1) - 1\bigr) = \tfrac12 + \tfrac16 = \tfrac23.
\]
The whole deck is worth 0 if you never start, and the full deck ends at zero if you do not stop, so all of the value lies in having the option to stop\index{value of the option to stop}. Note that the deck has a value even though every card is a fair bet. A prophet who stops at the top of the walk would make about 4 on a full deck, and the game gets about two-thirds of that with no foresight.

\iqwrong{Saying that the value is 0, by the optional stopping theorem. The running total is not a martingale: when you are ahead, more black cards than red remain, so the next card has a negative expected value. That is exactly when you stop.}
\end{solution}

\begin{solution}{iq:iv1:optimal-stopping:15}
Take the first free space that is at most 6 units before the restaurant, or failing that the first free space at or past it. The expected walk is about $6.57$. In general, with $q = 1 - p$, take a free space $x$ units before the restaurant exactly when $q^x \ge \tfrac12$.

Use backward induction\index{backward induction}. Past the restaurant, your walk only grows, so you take the first free space, which is $1/p$ units past on average. Let $V(x)$ be the expected walk on arriving $x$ units before the restaurant, before seeing that space. Then
\[
V(0) = \frac{q}{p}, \qquad V(x) = p\min\bigl(x, V(x - 1)\bigr) + q\,V(x - 1).
\]
If you take every free space from $x$ inwards, this solves to $V(x) = x + (2q^{x + 1} - q)/p$. At $x$ units, taking costs $x$, and passing costs $V(x - 1) = x - 1 + (2q^x - q)/p$. So you take the space when $p \le 2q^x - q$, which is $q^x \ge \tfrac12$\index{one-step look-ahead}. In words: take it when it is at least even money that all $x$ spaces from here up to and including the restaurant's are occupied.

With $q = 0.9$, we have $0.9^6 = 0.531$ and $0.9^7 = 0.478$, so the limit is 6 units, and $V(6) = 6 + (2 \cdot 0.9^7 - 0.9)/0.1 \approx 6.57$. Limits of 5 and 7 units give $6.63$ and $6.61$.

\iqwrong{Taking the first free space within 10 units, because the average gap between free spaces is $1/p = 10$. This costs $7.28$. The right test compares a chance with one half, not a distance with the average gap. Refusing everything before the restaurant costs $9$.}
\iqwhy{Once you would stop at some point, you would also stop at every point closer, so comparing stopping now with stopping at the next chance is enough. Rules of this monotone kind turn up in working an order before a deadline, and they save solving the whole recursion.}
\end{solution}

\begin{solution}{iq:iv1:optimal-stopping:16}
The profits are 0, $\tfrac13$, $\tfrac13$ and $\tfrac25$. The second free flip adds nothing, and the limit is $\tfrac12$. For odd $n$, the value is $\frac{n + 1}{2(n + 2)}$, and each even $n$ has the same value as the odd number before it.

With $k$ heads in $n$ flips, the chance of heads next is $\frac{k + 1}{n + 2}$, by Laplace's rule\index{Laplace's rule of succession}. You bet on the likelier side, so your profit is $|2\frac{k + 1}{n + 2} - 1|$. With a uniform prior, $k$ is uniform on $0, \dots, n$. So the value is the average of that absolute value over $k$.

For $n = 1$: chances $\tfrac13, \tfrac23$, and profit $\tfrac13$.

For $n = 2$: chances $\tfrac14, \tfrac12, \tfrac34$, and profits $\tfrac12, 0, \tfrac12$, so the mean is $\tfrac13$.

For $n = 3$: profits $\tfrac35, \tfrac15, \tfrac15, \tfrac35$, so the mean is $\tfrac25$.

Betting on the side of the first flip is still an optimal rule after the second flip. If the two agree, you bet that way anyway, and if they disagree, you are indifferent. So the second flip changes no decision, and it adds no value. The limit is $E|2p - 1| = \tfrac12$, which is what knowing $p$ exactly would be worth.

\iqwrong{Assuming that the value of data grows smoothly with $n$. It is flat here between an odd and an even step, because only a change in decision has value. Further data that sharpens your estimate without changing your bet adds nothing to a one-shot bet of fixed size.}
\end{solution}

\begin{solution}{iq:iv1:optimal-stopping:17}
A should keep a first draw of at least $1/\sqrt{3} \approx 0.577$, and A's chance of winning is about $0.462$.

Start with B\index{backward induction}. Given A's final number $x$, B wins unless both of B's draws are below $x$, because B redraws exactly when behind. So A wins with chance $x^2$.

Now, A has a draw $x$. Keeping gives $x^2$. Redrawing gives a uniform number, and so $E[U^2] = \tfrac13$. So A keeps when $x^2 \ge \tfrac13$. With threshold $t$, the chance that A wins is
\[
\int_t^1 x^2\,dx + t \cdot \tfrac13 = \frac{1 - t^3}{3} + \frac{t}{3},
\]
which is maximised at $t^2 = \tfrac13$, and has value $\tfrac13\bigl(1 + \tfrac{2}{3\sqrt{3}}\bigr) \approx 0.462$. As a check, $t = 0$ and $t = 1$ both mean that A has effectively one draw, which gives $\tfrac13$. A simulation of the game gives the same.

\iqwrong{Using $\tfrac12$ as the threshold, which is what maximises the expected number. A is not trying to get a high number on average; A has to survive B's two tries. The payoff is $x^2$, not $x$. Convexity favours a riskier choice.}
\end{solution}

\begin{solution}{iq:iv1:optimal-stopping:18}
Stop at the first success in the last 9 trials. The win chance is $0.9^9 \approx 0.387$. Hiring is the same problem with success chance $1/k$ at trial $k$.

Stopping at a success on trial $k$ wins if the remaining $n - k$ are all failures. So a success is worth more the later it comes, and the rule is a threshold on the trial number. For a rule that accepts any success in the last $s$ trials, you win when exactly one success occurs among them:
\[
P(\text{win}) = s p q^{s - 1},\quad q = 1 - p.
\]
Comparing $s$ with $s + 1$, the chance rises while $s \cdot p/q < 1$. Here $p/q = 1/9$, so $s = 9$ gives exactly 1, and $s = 9$ and $s = 10$ tie, at $0.387$. For $s = 8$ it is $0.383$.

This is the odds algorithm\index{odds algorithm}: sum the odds of success from the end backwards, and start accepting when the sum reaches 1. In the hiring problem, ``candidate $k$ beats all before'' is a success with chance $\frac1k$, these events are independent, and the last success is the overall best. The odds are $\frac{1}{k - 1}$. For $n = 10$, summing from the end reaches 1 at $k = 4$, so you reject the first 3, as found by direct calculation. The win chance is at least about $1/e$ when the sum of the odds reaches 1.

\iqwrong{Waiting until the expected number of remaining successes is about 1, which is the last 10 trials here. It is nearly right, and it gives the same win chance in this tied case. In general use the odds, not the probabilities, which matters when the chances are not small.}
\end{solution}

\begin{solution}{iq:iv1:optimal-stopping:19}
The rule ``stop at the first lead'' is worth exactly $\pi/4 \approx 0.785$. Backward induction with a long horizon gives about $0.793$ for the best rule, so that rule is not optimal. You should never stop when level or behind.

For the lower bound: a head on the first flip pays 1. After a first tail, the fair walk comes back to level with probability 1, so you can always reach a fraction of $\tfrac12$. This gives at least $\tfrac12 \cdot 1 + \tfrac12 \cdot \tfrac12 = \tfrac34$.

For the first lead\index{Chow--Robbins game}: the first lead comes at flip $2m + 1$ with chance $\frac{C_m}{2^{2m + 1}}$, with $C_m$ the Catalan numbers, and the fraction then is $\frac{m + 1}{2m + 1}$. So the value is
\[
\sum_{m \ge 0} \frac{C_m}{2^{2m + 1}} \cdot \frac{m + 1}{2m + 1} = \sum_{m \ge 0} \binom{2m}{m}\frac{1}{(2m + 1)2^{2m + 1}} = \frac12 \arcsin 1 = \frac{\pi}{4},
\]
using $C_m(m + 1) = \binom{2m}{m}$ and the arcsine series\index{arcsine series}.

Never stop when level or behind. From any state you can get back to at least $\tfrac12$ almost surely, because the walk is recurrent, so stopping below $\tfrac12$ loses. At exactly $\tfrac12$ you have a positive chance of getting a lead, and you can still come back, so continuing is strictly better.

The first lead is not always a good moment to stop either. After tail, head, head, you have $\tfrac23$, but going on is worth about $0.669$. So the optimal rule demands a bigger lead at some points, and the required lead grows like $0.84\sqrt{n}$. The value comes from a backward induction from 5000 flips, using a low and a high guess for the terminal value, and both round to $0.793$. It is the output of that calculation, not of a formula.

\iqwrong{Saying that the game is worth 1 because a fair walk gets ahead by any amount eventually. It gets ahead, but the fraction is what pays, and as the number of flips grows, any given lead is worth less. Being ahead by a fixed amount at a late time has a value close to $\tfrac12$.}
\end{solution}

\begin{solution}{iq:iv1:optimal-stopping:20}
Take the threshold $t = \tfrac12 E[M]$, with $M = \max X_i$, and accept the first prize that is at least $t$. That earns at least $\tfrac12 E[M]$. An example with one prize of 1 for sure, followed by one of $1/\varepsilon$ with chance $\varepsilon$, shows that you cannot do better than a half in general.

Let $p = P(M \ge t)$, the chance that you accept something\index{prophet inequality}. Your reward has two parts. Each time you accept you get at least $t$, which gives $pt$. On top of that, you get the excess $(X_i - t)^+$ when you accept prize $i$. You reach prize $i$ with no acceptance so far with chance at least $1 - p$, and that event depends only on the prizes before $i$, so it is independent of $X_i$. So
\[
\text{Reward} \ge pt + (1 - p)\sum_i E[(X_i - t)^+] \ge pt + (1 - p) E[(M - t)^+] \ge pt + (1 - p)(E[M] - t),
\]
where the sum is above $(M - t)^+$ because the largest excess is one of the terms. With $t = \tfrac12 E[M]$, both $t$ and $E[M] - t$ equal $\tfrac12 E[M]$, and $p$ cancels:
\[
\text{Reward} \ge p \cdot \tfrac12 E[M] + (1 - p) \cdot \tfrac12 E[M] = \tfrac12 E[M].
\]
For the example, the first prize is 1 for sure. The second prize is $1/\varepsilon$ with chance $\varepsilon$, so its mean is 1 too. So anything you do is worth 1. But the prophet takes the second prize when it comes, and otherwise the first, for $2 - \varepsilon$. The ratio tends to $\tfrac12$.

\iqwrong{Trying for a threshold that depends on all of the distributions in a complicated way. A single number already gets the best guarantee possible over all distributions, and the proof uses nothing about the order in which they arrive.}
\iqwhy{This gives the true value of hindsight in selection problems. Backtests that pick the best moment after the fact can show up to twice what any real rule can earn, even when the distributions are perfectly known.}
\end{solution}
